\chapter{Thảo luận}
\label{chap:discussion}

Kết quả thực nghiệm cho thấy bài toán phân bổ test-time compute không có một
điểm vận hành chung cho mọi tập dữ liệu. Dừng lấy mẫu theo đồng thuận tạo lợi
ích nhất quán khi so với cấu hình fixed $k=4$ tại cùng token cap. Ngược lại,
heuristic joint dựa trên độ dài prompt không chuyển giao ổn định từ GSM8K sang
MATH-500. Chương này diễn giải hai kết quả trên, sau đó phân loại các đe dọa độ
giá trị và xác định các đánh giá tiếp theo có ưu tiên cao nhất.

\section{Ý nghĩa của kết quả theo hai trục tài nguyên}

Hai trục $k$ và $L$ giải quyết hai dạng thiếu hụt khác nhau. Tăng $L$ cho một
quỹ đạo thêm cơ hội hoàn tất reasoning và xuất đáp án cuối; tăng $k$ tạo thêm
lần thử và cho phép tổng hợp đáp án. Khi token cap quá thấp, nhiều mẫu đều có
thể bị cắt trước đáp án cuối, nên breadth không bù được thiếu hụt về cap. Mẫu
hình này thể hiện rõ trên MATH-500: các cấu hình $L=1024$ có valid rate và
accuracy thấp dù tăng $k$.

Khi $L$ đã đủ lớn, lấy bốn mẫu cho mọi bài lại trở nên lãng phí. Trên GSM8K tại
$L=1024$, 54 câu đạt đồng thuận sau hai mẫu; trên MATH-500 tại $L=4096$, con số
này là 57 câu. Việc không mua hai mẫu cuối ở các câu này tạo mức giảm token
21.33\% và 19.96\% so với fixed $k=4$ tương ứng. Vector đúng--sai trong từng
cặp hoàn toàn giống nhau trên Test-100. Bằng chứng vì vậy hỗ trợ selective
breadth như một cơ chế tiết kiệm compute sau khi token cap đã ở mức phù hợp.

Kết quả không hàm ý rằng agreement là xác suất đúng đã được hiệu chuẩn. Hai
mẫu có thể đồng thuận trên cùng một đáp án sai. Trong hệ thống hiện tại,
agreement chỉ là điều kiện dừng chi phí. Giá trị của nó được đánh giá bằng kết
quả ghép cặp trên tập giữ lại, thay vì được giả định từ trước.

\section{Vì sao heuristic joint không ổn định}

Heuristic joint dùng tertile của độ dài prompt để chọn $L$, rồi dùng validity
và agreement để quyết định số mẫu. Thiết kế này có ưu điểm là rẻ và không dùng
đáp án chuẩn tại thời điểm quyết định, nhưng độ dài prompt là tín hiệu yếu cho
độ khó suy luận. Hai câu có độ dài gần nhau có thể khác lớn về số bước, kiến
thức toán và khả năng biểu diễn đáp án.

Trên GSM8K, routing theo độ dài vẫn tạo một điểm Pareto chi phí thấp: 70/100
với 1\,449.46 token. Tuy nhiên, trên MATH-500, joint chỉ đạt 36/100 với
5\,788.81 token và bị fixed $k=2,L=4096$ trội cả accuracy lẫn chi phí. Dữ liệu
trace cho thấy việc chọn cap thấp xảy ra trước; escalation theo validity hoặc
agreement ở cùng cap không thể phục hồi một quỹ đạo đã thiếu không gian sinh.
Sự khác biệt giữa hai tập bài cho thấy routing cần tín hiệu liên quan trực tiếp
hơn đến nhu cầu reasoning hoặc khả năng hoàn tất, thay vì chỉ dùng kích thước
đầu vào.

Oracle hồi cứu vẫn tìm thấy headroom vì được phép nhìn đáp án chuẩn và toàn bộ
pool của một seed. Headroom này không xác định được một rule triển khai. Một
mẫu đúng do may mắn có thể làm oracle lạc quan, trong khi policy thực không
biết trước mẫu nào sẽ đúng. Khoảng cách oracle do đó chỉ biểu diễn cơ hội trong
pool quan sát được; nó không phải ước lượng về accuracy có thể đạt của một
controller mới.

\section{Hàm ý cho mô hình suy luận cỡ nhỏ}

Với GPU cục bộ, compute cần được xem như một tài nguyên có giá thay vì chỉ là
tham số tăng chất lượng. Kết quả gợi ý một quy tắc triển khai thực dụng gồm hai
bước. Trước hết, chọn token cap đủ để valid rate không bị giới hạn nghiêm trọng
trên Validation của miền đích. Sau đó, dùng dừng breadth có thể kiểm toán để
tránh lấy đủ số mẫu tối đa cho mọi đầu vào.

Quy tắc trên không tạo ra một policy chung cho mọi miền. GSM8K chọn adaptive
breadth ở mức cap trung bình, còn MATH-500 chọn fixed $k=2$ tại cap cao làm
điểm vận hành chính. Vì vậy, đóng băng policy riêng theo dataset là lựa chọn
phù hợp hơn với bằng chứng hiện tại. Khi chuyển sang miền mới, cần lặp lại bước
Validation thay vì mang nguyên threshold và policy từ miền cũ.

Chi phí token là metric chính vì có thể so sánh giữa các lần replay tốt hơn
latency. Latency quan sát vẫn hữu ích cho lập kế hoạch chạy, nhưng phụ thuộc
GPU, trạng thái nhiệt, phiên bản phần mềm và cách batching. Một hệ thống triển
khai thực tế nên báo cáo cả token, latency và bộ nhớ, song không nên dùng một
phiên đo latency duy nhất làm kết luận tổng quát.

\section{Đe dọa độ giá trị cấu trúc}

Độ giá trị cấu trúc liên quan đến việc metric có đo đúng khái niệm nghiên cứu
hay không. Bốn giới hạn chính được xác định:

\begin{enumerate}
  \item $L$ là \code{max\_new\_tokens}, không phải độ sâu suy luận ngữ nghĩa.
  Một response dài hơn có thể lặp lại hoặc lan man mà không reasoning tốt hơn.
  \item Hard cap có thể cắt response trước đáp án cuối. Vì vậy, tác động của
  $L$ hiện trộn lẫn thêm thời gian suy luận với khả năng đơn thuần hoàn tất định
  dạng output.
  \item Exact-answer accuracy bỏ qua chất lượng reasoning trung gian. Parser
  nghiêm ngặt giúp metric tái lập nhưng có thể tính sai một response có ý đúng
  mà không tuân thủ định dạng đã đăng ký.
  \item Actual generated tokens đo chi phí sinh nhưng chưa đại diện đầy đủ cho
  năng lượng, bộ nhớ, batching hoặc thời gian người dùng chờ.
\end{enumerate}

Ba đối chứng decoding---budget forcing, native thinking budget của Qwen3 và
non-thinking mode---là cách trực tiếp nhất để tách hiệu ứng hoàn tất output
khỏi hiệu ứng deliberation. Chúng chưa được chạy, nên kết luận hiện tại chỉ áp
dụng cho thinking mode với hard token cap.

\section{Đe dọa độ giá trị nội tại}

Thiết kế shared pool và nested prefix giảm nhiễu do các policy nhận mẫu ngẫu
nhiên khác nhau. Manifest, model revision, prompt, parser, seed và policy ID
được khóa; replay kiểm tra pool thiếu, khóa trùng và generation error. Các biện
pháp này làm tăng khả năng quy khác biệt quan sát được cho policy trong cùng
pool.

Tuy vậy, nghiên cứu chỉ có một seed. Kết quả dừng breadth có vector đúng--sai
giống fixed $k=4$ trong pool hiện tại, nhưng pool ở seed khác có thể thay đổi
agreement và đáp án thắng. Ngoài ra, replay từ pool đã sinh mô phỏng quyết định
tuần tự về mặt mẫu và chi phí token; nó chưa tái tạo đầy đủ scheduler, batching
và latency của một dịch vụ online thật.

Validation chỉ có 20 câu mỗi dataset, nên lựa chọn policy có độ nhiễu cao. Đối
với MATH-500, pivot sang fixed $k=2,L=4096$ được thực hiện và ghi nhận sau
Validation nhưng trước Test. Trình tự này bảo vệ Test khỏi tuning, song không
làm cho ước lượng Validation trở nên chính xác hơn. Một split Validation lớn
hơn hoặc nested resampling sẽ giảm phụ thuộc vào 20 câu cụ thể.

\section{Đe dọa độ giá trị ngoại tại}

Bằng chứng chính chỉ dùng Qwen3-0.6B, một model revision và hai benchmark toán.
GSM8K và MATH-500 khác nhau về độ khó nhưng chưa đại diện cho coding, lập luận
logic, khoa học hoặc ngôn ngữ tự nhiên nói chung. Kết quả cũng chưa cho biết
policy có chuyển giao sang mô hình nhỏ của họ khác hay sang kích thước lớn hơn
hay không.

Pilot Qwen3-1.7B trên GSM8K hỗ trợ việc hình thành giả thuyết nhưng không phải
đối chứng ngoại tại hoàn chỉnh. Thực nghiệm 1.7B trên MATH-500 không hoàn tất
do giới hạn thời gian và bộ nhớ của RTX 4060. Do đó, luận văn chỉ khái quát kết
quả đến cấu hình 0.6B đã kiểm tra; mọi phát biểu rộng hơn được giữ ở mức giả
thuyết.

\section{Đe dọa độ giá trị kết luận}

Mỗi Test có 100 câu. Khoảng Wilson biểu diễn bất định do lấy mẫu câu hỏi, còn
McNemar exact khai thác ghép cặp đúng--sai giữa hai policy. Hiệu chỉnh Holm làm
giảm nguy cơ dương tính giả trong họ so sánh xác nhận. Tuy nhiên, các khoảng và
kiểm định này không bao phủ biến thiên giữa seed, model revision hoặc lựa chọn
Validation.

Pareto membership được tính trên 14 policy đã đăng ký và mang tính mô tả. Nó
có thể thay đổi khi bổ sung policy mới. Comparator breadth $L=4096$ trên
MATH-500 là kết quả phụ đã đăng ký; accuracy 73/100 của nó không thay thế ước
lượng chính 66/100 từ policy selected. Việc giữ nguyên hai vai trò này tránh
chọn lại policy theo Test.

\begin{table}[htbp]
  \centering
  \caption{Tóm tắt đe dọa độ giá trị và biện pháp kiểm soát.}
  \label{tab:validity-threats}
  \scriptsize
  \begin{tabular}{p{0.16\textwidth}p{0.31\textwidth}p{0.39\textwidth}}
    \toprule
    Nhóm & Đe dọa chính & Kiểm soát hoặc ranh giới diễn giải \\
    \midrule
    Cấu trúc & Token cap không đồng nghĩa semantic depth; hard truncation. &
    Báo cáo valid/truncation và giới hạn claim ở hard-cap thinking mode. \\
    Nội tại & Một seed; Validation nhỏ; replay chưa phải serving online. &
    Shared pool, freeze record, paired analysis và mô tả pivot trước Test. \\
    Ngoại tại & Một model 0.6B và hai benchmark toán. &
    Không gộp pilot 1.7B; không khái quát sang họ mô hình hoặc miền khác. \\
    Kết luận & Test-100, nhiều policy và Pareto phụ thuộc registry. &
    Wilson, McNemar exact, Holm; Pareto chỉ mô tả và không dùng để chọn lại. \\
    \bottomrule
  \end{tabular}
\end{table}

\section{Đánh giá tiếp theo theo thứ tự ưu tiên}

Trong phạm vi thời gian còn lại, công việc mới chỉ nên được thêm nếu không làm
thay đổi đánh giá Test chính đã đóng băng. Thứ tự ưu tiên là:

\begin{enumerate}
  \item chạy ba decoding control trên Validation hoặc một split robustness
  được khai báo riêng, tập trung vào các điểm vận hành đã chọn thay vì mở lại
  toàn bộ lưới;
  \item lặp thêm seed cho selected policy và comparator trực tiếp để ước lượng
  độ nhạy ngẫu nhiên với chi phí GPU kiểm soát được;
  \item dùng phần reserve MATH-500 làm confirmatory set mới nếu bất kỳ threshold
  hoặc controller nào tiếp tục được phát triển;
  \item chỉ huấn luyện learned controller sau khi thu thập đủ escalation event,
  sau đó đánh giá trên một Test chưa từng dùng;
  \item kiểm tra chuyển giao trên một mô hình nhỏ khác và một miền ngoài toán.
\end{enumerate}

Ba mục đầu có giá trị cao nhất cho độ tin cậy của kết luận hiện tại. Hai mục
sau là hướng mở rộng nghiên cứu, không phải điều kiện bắt buộc để hoàn thành
khóa luận ở quy mô hiện nay.

